% ============================================================================
% SECCIÓN: JUSTIFICACIÓN DE LOS FUNDAMENTOS TEÓRICOS
% ============================================================================

\subsection{Justificación de los fundamentos teóricos seleccionados}

SocialScrum enfrenta tres preguntas fundamentales sobre los equipos ágiles que no pueden responderse únicamente con perspectivas técnicas:

\begin{enumerate}
\item ¿Qué motiva realmente a las personas en equipos de desarrollo de software, especialmente cuando la presión aumenta?
\item ¿Qué mecanismos sostienen la colaboración efectiva en entornos complejos y mediados por tecnología?
\item ¿Cómo prevenir el deterioro social que genera deuda y compromete el bienestar?
\end{enumerate}

La selección de la Teoría de la Autodeterminación (TAD) y el Modelo Integrador de Confianza (Mayer et al.) responde directamente a estas preguntas y refleja la evidencia empírica sobre qué impide o favorece la deuda social.

\subsubsection{¿Por qué la Teoría de la Autodeterminación (TAD) y no otras teorías motivacionales?}

Existen múltiples teorías sobre motivación humana en contextos laborales: la Teoría de Expectativas (Vroom), la Teoría de Metas (Locke y Latham), la Teoría de dos Factores (Herzberg), y la Teoría de Necesidades (McClelland). Sin embargo, TAD es superior para abordar la deuda social por tres razones clave.

\paragraph{Razón 1: Enfoque en motivación intrínseca versus síntomas de burnout}

La Teoría de Expectativas predice que la motivación depende del valor esperado de una recompensa multiplicado por la probabilidad de obtenerla. La Teoría de Metas enfatiza el establecimiento claro de objetivos. Ambas se centran en mecanismos extrínsecos: dinero, reconocimiento, plazos.

Sin embargo, la deuda social y el burnout no se previenen únicamente con recompensas externas. De hecho, investigaciones recientes muestran que una dependencia excesiva de incentivos externos debilita la motivación intrínseca \cite{ryan_self-determination_2000, tulili2022burnout}. El burnout surge precisamente cuando las personas pierden el sentido de autonomía, sienten que sus esfuerzos no producen competencia (ineficacia) y se desconectan de sus compañeros (falta de relacionalidad) \cite{tulili2022burnout}.

TAD, por el contrario, propone que existen tres necesidades psicológicas universales cuya satisfacción sostiene la motivación intrínseca: autonomía, competencia y relacionalidad. Estas tres necesidades mapean directamente con los tres componentes del síndrome de burnout (agotamiento, ineficacia, cinismo) \cite{ryan_self-determination_2000, tulili2022burnout}.

\textbf{Conclusión:} TAD explica no solo cómo motivar, sino cómo prevenir el colapso motivacional que caracteriza a la deuda social. Sin ella, SocialScrum solo vería síntomas, no causas.

\paragraph{Razón 2: Universalidad transcultural y aplicabilidad en equipos diversos}

TAD ha sido validada en más de 100 estudios en contextos culturales variados: desde individuos en occidente hasta equipos en Asia, América Latina y África \cite{ryan_self-determination_2000}. Esta robustez es esencial para un marco que se pretende aplicable en entornos globales y diversos.

En contraste, la Teoría de Metas presenta variaciones significativas según contexto cultural (culturas individualistas vs. colectivistas), mientras que la Teoría de dos Factores de Herzberg ha sido cuestionada por su estrecha metodología y falta de validación en contextos actuales.

\textbf{Conclusión:} TAD puede ser adaptada a diferentes culturas sin perder su núcleo explicativo. Esto es crítico para un modelo ágil global.

\paragraph{Razón 3: Conexión empírica con community smells}

Estudios sobre community smells específicos revelan que los patrones de aislamiento (lone wolf), falta de cooperación (sharing villainy), y burn-out correlacionan directamente con déficits en las tres necesidades de TAD \cite{caballero2022communitysmells, saeeda2024nontechnicaldebt}.

Por ejemplo, el community smell ``lone wolf'' (lobo solitario) emerge cuando los desarrolladores perciben falta de autonomía, no confían en su competencia relativa al equipo, o se sienten desconectados emocionalmente. TAD ofrece un diagnóstico preciso: ¿cuál de las tres necesidades se ve comprometida?

Ninguna otra teoría motivacional proporciona este nivel de especificidad diagnóstica para community smells.

\textbf{Conclusión:} TAD no es solo un marco motivacional, es un diagrama que conecta síntomas observables (community smells) con causas psicológicas profundas.

\subsubsection{¿Por qué el Modelo de Confianza de Mayer et al. y no otros modelos de trabajo en equipo?}

Existen modelos alternativos para entender dinámicas de equipo: el Modelo de Etapas de Tuckman (formación, tormenta, normalización, desempeño), el Modelo de Psicología Grupal de Bion, o la Teoría del Capital Social de Putnam. ¿Por qué Mayer et al.?

\paragraph{Razón 1: La confianza es el mecanismo generativo de la deuda social}

Tuckman explica cómo evolucionan los equipos, pero no explica qué causa que una etapa se estanque o deteriore. Bion se enfoca en dinámicas inconscientes, pero su lenguaje es menos pragmático para un equipo ágil.

Sin embargo, la literatura sobre deuda social es clara: la deuda social crece en entornos donde la confianza colapsa \cite{wilson_all_2006, mayer_integrative_1995}. Cuando la confianza falla:
\begin{itemize}
\item Aumenta el monitoreo mutuo (que genera ineficiencia)
\item Se fragmenta la comunicación (aparece radio silence)
\item Se duplican tareas (knowledge hoarding)
\item Surge aislamiento (lone wolf)
\end{itemize}

El Modelo de Confianza de Mayer et al. propone tres factores operacionales: habilidad, benevolencia e integridad. Estos factores son observables, medibles y accionables. A diferencia de otros modelos más descriptivos, este es prescriptivo: si diagnosticas que falla la benevolencia, sabes qué intervenir.

\textbf{Conclusión:} La confianza no es un epifenómeno de equipos saludables; es el mecanismo que previene o genera deuda social. Mayer et al. proporciona el mapa para navegar ese mecanismo.

\paragraph{Razón 2: Operacionalización clara de los factores}

El Capital Social de Putnam es conceptualmente elegante, pero difícil de operacionalizar en un equipo de desarrollo concreto. ¿Cómo midió exactamente la falta de capital social en su estudio?

Mayer et al., por el contrario, ofrece criterios concretos:
\begin{itemize}
\item \textbf{Habilidad:} ¿Confían los miembros en las competencias técnicas y sociales de cada uno?
\item \textbf{Benevolencia:} ¿Perciben que sus compañeros actúan con genuina preocupación por su bienestar?
\item \textbf{Integridad:} ¿Observan coherencia entre lo que se dice y lo que se hace?
\end{itemize}

Estos criterios pueden evaluarse mediante surveys, entrevistas o análisis de red social (SNA). Son traducibles a métricas.

\textbf{Conclusión:} SocialScrum necesita un modelo operacionable. Mayer et al. lo proporciona.

\paragraph{Razón 3: Validación empírica en contextos técnicos}

El Modelo de Confianza de Mayer ha sido validado específicamente en equipos de trabajo distribuidos y mediados por tecnología \cite{wilson_all_2006, mayer_integrative_1995}. Este es exactamente el contexto de los equipos ágiles modernos.

En contraste, el Modelo de Etapas de Tuckman fue desarrollado en los años 1960s en contextos principalmente presenciales y no se ha validado completamente en equipos remotos o híbridos.

\textbf{Conclusión:} Mayer et al. está calibrado para la realidad actual: equipos remotos, mediados y bajo presión.

\subsubsection{¿Qué NO puedes explicar sin estas teorías?}

Sin TAD:
\begin{itemize}
\item No puedes diferenciar entre falta de motivación temporal (agotamiento) y colapso motivacional crónico (síndrome de burnout)
\item No puedes predecir por qué un equipo con todas las metas cumplidas sigue siendo tóxico
\item No puedes intervenir estratégicamente: una bonificación (refuerzo externo) puede empeorar la motivación intrínseca
\end{itemize}

Sin el Modelo de Confianza:
\begin{itemize}
\item No puedes diagnosticar específicamente dónde falla la confianza (¿falta habilidad? ¿benevolencia? ¿integridad?)
\item No puedes distinguir entre equipos disfuncionales por competencia vs. por falta de empatía
\item No tienes un criterio claro para cuándo la intervención debe ser técnica (capacitación) vs. relacional (mediación)
\end{itemize}

% ============================================================================
% SECCIÓN: EXPANSIÓN DE VALORES DE SCRUM CON COMMUNITY SMELLS
% ============================================================================

\subsection{Aplicación extendida de los valores de Scrum en SocialScrum con mitigación de community smells}

Cada valor de Scrum es un mecanismo directo de prevención o mitigación de community smells específicos. A continuación, se desarrolla cómo cada valor opera en el contexto de SocialScrum para abordar los patrones disfuncionales observables en equipos de desarrollo.

\paragraph{Coraje (Courage)}

El coraje es el fundamento emocional que permite al equipo actuar sobre la vulnerabilidad. La confianza, como se ha definido, requiere precisamente disposición a ser vulnerable ante otros \cite{mayer_integrative_1995}. Sin coraje, esa vulnerabilidad se percibe como riesgo, no como oportunidad.

\subsubsection*{Definición profunda}

El coraje en SocialScrum no es la ausencia de miedo, sino la capacidad de actuar a pesar del miedo a consecuencias sociales negativas. Se expresa a través del concepto de \textit{Risk Taking in Relationship} (RTR): la disposición a expresar desacuerdo, pedir ayuda, admitir errores, o confrontar dinámicas tóxicas sin temor a represalias sociales o exclusión del grupo.

\subsubsection*{Community smells prevenidos o mitigados por Coraje}

\begin{itemize}

\item \textbf{Radio Silence:} Este smell emerge cuando los miembros evitan comunicar problemas, conflictos o preocupaciones. Sin coraje, los desarrolladores observan un problema técnico-social (por ejemplo, una mala decisión de arquitectura impulsada por presión) pero optan por el silencio para evitar confrontación. El coraje rompe ese silencio permitiendo que los problemas salgan a la luz tempranamente \cite{saeeda2024nontechnicaldebt}.

\item \textbf{Sharing Villainy:} Este smell describe la falta de voluntad para compartir conocimiento crítico, a menudo por miedo a ser reemplazado o porque se percibe que el conocimiento es poder. El coraje permite que los desarrolladores superen esa inseguridad y confíen en que contribuir al bien común refuerza (no debilita) su posición en el equipo \cite{caballero2022communitysmells}.

\item \textbf{Lone Wolf (Lobo Solitario):} Este smell refleja aislamiento, donde un miembro trabaja sin colaboración. Frecuentemente surge porque la persona tiene miedo a no ser competente con sus pares o temor al rechazo. El coraje permite que se acerque pesar del riesgo de vulnerabilidad \cite{caballero2022communitysmells}.

\item \textbf{Lack of Conflict Resolution:} Sin coraje, los conflictos interpersonales se ignoran, se reprimen o se manejan pasivo-agresivamente. El coraje permite que el equipo y el Social Guide aborden los conflictos de forma directa y constructiva, utilizando la \textit{Sprint Retrospective Social} como un espacio seguro para esa conversación \cite{saeeda2024nontechnicaldebt}.

\end{itemize}

\subsubsection*{Mecanismo de acción en SocialScrum}

El Social Guide crea espacios explícitos de seguridad psicológica que disminuyen la percepción de riesgo social. Estrategias concretas incluyen:
\begin{enumerate}
\item \textbf{Normas de grupo explícitas:} Al inicio del Sprint, se establecen acuerdos sobre cómo se manejarán los desacuerdos y las críticas (ej.: ``las críticas son sobre ideas, no sobre personas'')
\item \textbf{Validación emocional:} Cuando alguien expresa preocupación o conflicto, se valida el valor de su contribución, no se castiga
\item \textbf{Rituales de disculpa y reparación:} Si alguien comete error social (ej.: lenguaje descortés), el equipo cuenta con un ritual breve de reconocimiento y reparación, demostrando que los errores sociales son reparables
\end{enumerate}

\subsubsection*{Evidencia empírica}

El caso Boeing (reporte de 2024) documenta cómo la falta de coraje en comunicar problemas de seguridad derivó en deuda social catastrófica: empleados temían reportar defectos de manufactura por miedo a represalias, lo que generó acumulación de riesgo técnico y colapso de confianza organizacional \cite{saeeda2024nontechnicaldebt}. Este es un ejemplo extremo, pero ilustra cómo Radio Silence y Lack of Conflict Resolution operan en sistemas sin coraje.

\paragraph{Apertura (Openness)}

La apertura es inseparable del pilar empírico de transparencia. Sin embargo, va más allá de simplemente visibilizar trabajo técnico; se trata de visibilizar la dimensión social, emocional y relacional del equipo.

\subsubsection*{Definición profunda}

La apertura en SocialScrum significa revelación voluntaria y honesta de:
\begin{itemize}
\item Percepciones sobre la dinámica del equipo
\item Dificultades personales que afectan el trabajo (sin intrusión en privacidad)
\item Preocupaciones sobre decisiones socio-técnicas
\item Estado emocional y necesidades de apoyo
\end{itemize}

Esta revelación solo ocurre cuando existe confianza psicológica; a su vez, la apertura mutua refuerza la confianza \cite{wilson_all_2006}.

\subsubsection*{Community smells prevenidos o mitigados por Apertura}

\begin{itemize}

\item \textbf{Radio Silence:} Al hacer explícitos los canales de comunicación y la legitimidad de hablar sobre problemas, la apertura reduce dramáticamente el silencio. El Social Guide puede preguntar directamente: ``¿Hay algo que no estamos diciendo?'' \cite{saeeda2024nontechnicaldebt}.

\item \textbf{Lone Wolf:} El aislamiento persiste porque es invisible. La apertura requiere que el aislado (o sus compañeros) verbalice lo que está pasando. Una vez visible, puede abordarse \cite{caballero2022communitysmells}.

\item \textbf{Organizational Silo:} Este smell describe silos de conocimiento entre subequipos. La apertura sobre quién sabe qué, dónde existen brechas, y cuáles son las dependencias permite romper los silos deliberadamente \cite{caballero2022communitysmells}.

\item \textbf{Lack of Information Sharing:} Directamente relacionado. La apertura obliga a que la información fluya, no se hoarde.

\end{itemize}

\subsubsection*{Mecanismo de acción en SocialScrum}

Estrategias operacionales:
\begin{enumerate}
\item \textbf{Daily Social Check-in:} Pregunta explícita en el Daily Scrum: ``¿Cómo te sientes hoy con el equipo? ¿Hay algo que afecte la colaboración?''
\item \textbf{Social Metrics Visible:} Se publican métricas sobre salud del equipo (índice de confianza, mapa de red social, prevalencia de smells) en la velocidad que el equipo se siente cómodo compartiendo
\item \textbf{Retrospective openness protocol:} Protocolo estructurado donde cada miembro habla brevemente sobre qué observó que benefició/daño la dinámica
\end{enumerate}

\paragraph{Respeto (Respect)}

El respeto se conecta con la necesidad de relacionalidad de TAD y con la benevolencia del modelo de confianza. Es el fundamento de una comunidad verdaderamente segura.

\subsubsection*{Definición profunda}

El respeto es el reconocimiento activo de la dignidad, competencia y aportación única de cada miembro. No es tolerancia pasiva (permitir que alguien esté en el equipo sin molestar), sino reconocimiento activo de valor.

\subsubsection*{Community smells prevenidos o mitigados por Respeto}

\begin{itemize}

\item \textbf{Toxic Relations:} Este smell describe relaciones hostiles, lenguaje descortés, o descalificación personal. Es lo opuesto directo del respeto. Un equipo que practica respeto activo no tolera la descortesía; la interviene inmediatamente \cite{caballero2022communitysmells}.

\item \textbf{Burnout:} El síndrome de burnout incluye cinismo (despersonalización). El cinismo emerge cuando una persona se siente desrespetada, sus aportaciones no son reconocidas, y percibe que sus sentimientos no importan. El respeto activo es un antídoto directo \cite{tulili2022burnout}.

\item \textbf{Lone Wolf:} A menudo, el aislamiento es resultado de falta de respeto previo. Si una persona fue descalificada o desestimada en el pasado, se aísla como mecanismo de defensa. El respeto crea las condiciones para que se reintegre.

\item \textbf{Lack of Conflict Resolution:} Los conflictos sin respeto mútuo se vuelven personales y se perpetúan. El conflicto respetuoso es aquel donde ambas partes sienten que su posición es escuchada, incluso si no están de acuerdo \cite{saeeda2024nontechnicaldebt}.

\end{itemize}

\subsubsection*{Mecanismo de acción en SocialScrum}

Estrategias operacionales:
\begin{enumerate}
\item \textbf{Reconocimiento de contribuciones:} En cada Sprint Review y Retrospective, el Social Guide destaca contribuciones específicas de cada miembro, incluyendo aportaciones técnicas y sociales
\item \textbf{Norm enforcement:} Si alguien es descortés (lenguaje descalificador, interrupción frecuente, etc.), se interviene en el momento y se alinea con el valor de respeto del equipo
\item \textbf{Vulnerability leadership:} El Social Guide y Scrum Master modelan el respeto siendo vulnerables primero: reconociendo sus propios errores, admitiendo no saber algo, pidiendo ayuda. Esto da permiso al equipo para hacer lo mismo
\end{enumerate}

\paragraph{Compromiso (Commitment)}

El compromiso en SocialScrum va más allá del Sprint Goal. Es el compromiso de cada miembro con el bienestar colectivo, incluso cuando esto requiere esfuerzo adicional.

\subsubsection*{Definición profunda}

El compromiso es la disposición sostenida a contribuir al bienestar del equipo, a asumir responsabilidad sobre problemas que afectan la dinámica, y a participar activamente en la prevención y mitigación de deuda social.

\subsubsection*{Community smells prevenidos o mitigados por Compromiso}

\begin{itemize}

\item \textbf{Organizational Silo:} Si hay compromiso verdadero con el equipo como entidad unificada, los silos de conocimiento se rompen deliberadamente. La persona comprometida pregunta: ``¿Qué información que tengo, otros necesitan?'' \cite{caballero2022communitysmells}.

\item \textbf{Sharing Villainy:} El compromiso requiere acción explícita para compartir conocimiento y recursos, incluso si eso significa reducir el control monopolístico sobre expertise \cite{caballero2022communitysmells}.

\item \textbf{Lack of Ownership:} Este smell describe equipos donde nadie siente responsabilidad real por los resultados. El compromiso genuino transforma esa dinámica: cada miembro siente que el resultado es en parte responsabilidad suya \cite{saeeda2024nontechnicaldebt}.

\item \textbf{Turnover y Team Instability:} Indirectamente, el compromiso con el equipo reduce la rotación no deseada. Si los miembros sienten pertenencia e importancia, tienden a permanecer y contribuir.

\end{itemize}

\subsubsection*{Mecanismo de acción en SocialScrum}

Estrategias operacionales:
\begin{enumerate}
\item \textbf{Shared Sprint Goal Social:} Además del Sprint Goal técnico, el equipo define un objetivo social para el Sprint (ej.: ``Reducir Radio Silence mediante check-ins diarios más auténticos'')
\item \textbf{Collective accountability:} Si un community smell persiste, se trata no como problema de una persona, sino como responsabilidad colectiva. ``¿Qué estamos haciendo colectivamente que permite que Lone Wolf persista?''
\item \textbf{Role rotation:} El Social Guide puede rotar entre miembros, demostrando que la responsabilidad por la salud del equipo es compartida, no delegada
\end{enumerate}

\paragraph{Foco (Focus)}

El foco es mantener la atención en lo que realmente importa. En SocialScrum, esto significa priorizar tanto la entrega técnica como la salud social, y resolver las tensiones cuando compiten.

\subsubsection*{Definición profunda}

El foco es la capacidad sostenida de priorizar lo esencial y resistir la fragmentación de atención. En contexto de deuda social, significa elegir conscientemente qué community smells abordar primero, basado en impacto e urgencia.

\subsubsection*{Community smells prevenidos o mitigados por Foco}

\begin{itemize}

\item \textbf{Burnout:} El agotamiento frecuentemente resulta de intentar todo a la vez. El foco permite que el equipo reconozca que no pueden resolver todos los problemas de una vez; deben ser estratégicos. Esto reduce la sensación de sobrecarga \cite{tulili2022burnout}.

\item \textbf{Ping-Pong Effect:} Este smell, aunque menos documentado, describe equipos que saltan entre prioridades sin resolver nada profundamente. El foco requiere que se aborde un community smell a la vez, hasta que sea resuelto o mitigado \cite{caballero2022communitysmells}.

\item \textbf{Lack of Conflict Resolution:} Cuando hay foco compartido en resolver un conflicto específico, se dedican recursos (tiempo, energía emocional, facilitation). Sin foco, el conflicto se ignora perpetuamente.

\end{itemize}

\subsubsection*{Mecanismo de acción en SocialScrum}

Estrategias operacionales:
\begin{enumerate}
\item \textbf{Smell prioritization:} En cada Sprint Retrospective Social, el equipo identifica todos los smells pero deliberadamente elige 1-2 para trabajar en el próximo Sprint. El resto se documentan pero se posponen \cite{saeeda2024nontechnicaldebt}.
\item \textbf{Social Backlog prioritization:} El Product Backlog Social (elementos de deuda social a abordar) es priorizado igual que el Product Backlog técnico. El foco requiere que ambos sean visibles en la priorización.
\item \textbf{Protected time:} Se asigna tiempo explícito en el Sprint para trabajar en acciones de mitigación social. Este tiempo no es negociable, incluso bajo presión de entrega técnica.
\end{enumerate}

% ============================================================================
% SECCIÓN: PRINCIPIOS EMPÍRICOS APLICADOS A DEUDA SOCIAL
% ============================================================================

\subsection{Aplicación de los pilares empíricos de Scrum a la deuda social}

Scrum se fundamenta en tres pilares empíricos: transparencia, inspección y adaptación \cite{schwaber2020scrum}. Estos pilares son aplicables no solo a la entrega técnica, sino a la gestión integral de la deuda social. A continuación, se desarrolla cómo cada pilar operacionaliza la prevención y mitigación de community smells.

\subsubsection{Transparencia aplicada a la deuda social}

\paragraph{¿Por qué es crítica para la deuda social específicamente?}

La deuda social es por naturaleza invisible. Los community smells no aparecen en burndown charts, no generan excepciones en logs, y no son detectados por herramientas de código estático. Sin mecanismos explícitos de visibilización, permanecen ocultos hasta que causan daño significativo: rotación de personal, declive de productividad, conflictos abiertos, o colapso de confianza \cite{saeeda2024nontechnicaldebt, caballero2022communitysmells}.

La deuda técnica es visible (hay código que incumple estándares, hay tests que fallan). La deuda social permanece oculta en las dinámicas humanas, en patrones de comunicación, en lo no dicho. Sin transparencia deliberada, el equipo ignora que existe.

\paragraph{Consecuencias documentadas de perder transparencia en deuda social}

\begin{itemize}
\item \textbf{Acumulación silenciosa:} Sin visibilización, la deuda social se acumula exponencialmente. Lo que empieza como una falta de comunicación menor se convierte en desconfianza sistémica en meses \cite{tamburri2015insights}.
\item \textbf{Diagnóstico tardío:} Cuando un problema se hace evidente (ej.: dimisión de miembro clave), ya es tarde para prevención; solo queda remediación costosa.
\item \textbf{Ciclos de represalia:} Sin visibilidad, los conflictos se manejan de forma implícita o pasivo-agresiva, generando ciclos donde cada parte se siente víctima de la otra.
\end{itemize}

El caso de Uber (2017): la investigación post-crisis reveló que existían claros indicadores de cultura tóxica que no eran visibles formalmente en reportes. Los smells (acoso, sexismo, falta de seguridad psicológica) permanecieron invisibles hasta que salieron a luz públicamente, causando daño masivo a la empresa y sus empleados \cite{saeeda2024nontechnicaldebt}.

\paragraph{Mecanismos de transparencia específicos en SocialScrum}

\begin{enumerate}

\item \textbf{Product Backlog Social:} Artefacto que registra explícitamente los elementos de deuda social identificados. A diferencia del trabajo técnico que aparece en el Product Backlog regular, este backlog es dedicado:
\begin{itemize}
\item ``Reducir Radio Silence mediante protocolo de comunicación asincrónica estructurada''
\item ``Mitigar Lone Wolf en [nombre del desarrollador] mediante pair programming deliberado''
\item ``Resolver conflicto de confianza entre [roles] mediante mediación estructurada''
\end{itemize}

Estos ítems son visibles, estimados, y priorizados como cualquier otra tarea. Su visibilidad reconoce que la deuda social requiere inversión y tiempo, igual que la deuda técnica \cite{saeeda2024nontechnicaldebt}.

\item \textbf{Social Metrics Dashboard:} Conjunto de indicadores que miden la salud social del equipo:
\begin{itemize}
\item \textbf{Índice de confianza:} Medido mediante surveys breves (escala 1-5) después de cada Sprint sobre si los miembros confían en las competencias, intenciones e integridad de sus compañeros
\item \textbf{Community Smell Registry:} Conteo de smells identificados por Sprint, con seguimiento de si se están resolviendo o acumulando
\item \textbf{Comunicación Asincrónica:} Métrica de cuántas conversaciones importantes ocurren en canales públicos vs. privados (Low private/high public = transparencia mayor)
\item \textbf{Velocidad de Resolución de Conflictos:} Tiempo promedio entre identificación de un conflicto y su resolución
\end{itemize}

Estos indicadores no reemplazan medidas de desempeño técnico, pero complementan el panorama. Un equipo que entrega rápido pero tiene conflictos no resueltos está acumulando deuda social.

\item \textbf{Sprint Retrospective Social:} Ceremonia donde la revisión de transparencia es central. El Social Guide presenta:
\begin{itemize}
\item Lista de smells identificados (nuevos y persistentes)
\item Cambios en métricas de confianza vs. Sprint anterior
\item Análisis de red social (si se utiliza): quién está aislado, cuáles son los puentes de información, dónde hay silos
\end{itemize}

Esta visibilización crea accountability: ``Vemos que Radio Silence persiste. ¿Qué vamos a hacer diferente?''

\item \textbf{Comunicación Explícita del Impacto:} El Social Guide conecta la deuda social con resultados concretos. No es suficiente decir ``hay un problema de confianza''. Es necesario explicitar: ``Cuando falló la confianza en [situación específica], resultó en [decisión técnica subóptima] que nos costó [cantidad] de tiempo de retrabazo'' \cite{saeeda2024nontechnicaldebt}.

\end{enumerate}

\subsubsection{Inspección aplicada a la deuda social}

\paragraph{¿Por qué es crítica para la deuda social específicamente?}

Hacer visible un problema no es suficiente; se debe inspeccionar regularmente para entender sus causas, su evolución, y su impacto. La inspección en contexto de deuda social significa análisis sistemático de patrones de interacción, dinámicas de poder, y mecanismos de deterioro social.

Sin inspección, la visibilización es superficial: el equipo ve que existe un smell, pero no entiende por qué, y por tanto no puede intervenir efectivamente \cite{tamburri2015insights}.

\paragraph{Dimensiones de la inspección de deuda social}

La inspección debe operar en tres niveles simultáneamente:

\begin{enumerate}

\item \textbf{Inspección de síntomas (nivel observable):} ¿Cuáles son los patrones visibles?
\begin{itemize}
\item ¿Cuántas conversaciones técnicas importantes ocurren en privado en lugar de en canales públicos?
\item ¿Quién inteviene regularmente en discussions y quién se queda en silencio?
\item ¿Cuántos conflictos se mencionan explícitamente vs. cuántos permanecen implícitos?
\end{itemize}

\item \textbf{Inspección de causas (nivel explicativo, conectado a TAD/Confianza):} ¿Por qué existen estos patrones?
\begin{itemize}
\item ¿Faltan las necesidades de autonomía, competencia, o relacionalidad? (TAD)
\item ¿La desconfianza surge de falta de habilidad percibida, dudas sobre benevolencia, o inconsistencia observada? (Confianza)
\item ¿El community smell mapea a uno o más de estos déficits fundamentales?
\end{itemize}

\item \textbf{Inspección de impacto (nivel sistémico):} ¿Cómo afecta la deuda social los resultados del equipo?
\begin{itemize}
\item ¿La baja confianza en competencia resultó en duplicación de trabajo?
\item ¿La falta de comunicación produjo decisiones técnicas subóptimas?
\item ¿El burnout correlaciona con períodos de mayor aislamiento social?
\end{itemize}

\end{enumerate}

\paragraph{Mecanismos de inspección específicos en SocialScrum}

\begin{enumerate}

\item \textbf{Social Network Analysis (SNA):} Análisis de redes sociales aplicado al equipo:
\begin{itemize}
\item Se mapea quién se comunica con quién, sobre qué temas, y con qué frecuencia
\item Se identifican hubs (personas centrales en la comunicación), bridges (personas que conectan grupos), y isolates (personas con bajo grado de conexión)
\item El análisis se repite cada 2-4 Sprints para detectar cambios de patrón \cite{saeeda2024nontechnicaldebt, caballero2022communitysmells}
\end{itemize}

SNA no requiere software sofisticado; puede hacerse con herramientas simples o incluso manualmente en equipos pequeños.

\item \textbf{Root Cause Analysis (RCA) de Community Smells:} Para cada smell identificado, se aplica análisis de causa raíz:
\begin{itemize}
\item \textbf{Síntoma}: ``Radio Silence: [nombre] no comunica problemas durante Daily Scrum''
\item \textbf{Causas posibles} (niveles):
  \begin{itemize}
  \item Nivel 1 (observable): No sabe cómo reportar el problema, le falta la información necesaria
  \item Nivel 2 (interpersonal): Teme ser juzgado o descalificado por sus compañeros
  \item Nivel 3 (sistémico): Experiencias previas de represalia cuando comunica problemas
  \item Nivel 4 (psicológico/teórico): Déficit de autonomía (no cree que sus preocupaciones sean válidas) o déficit de relacionalidad (no se siente seguro en el grupo)
  \end{itemize}
\item \textbf{Hipótesis:} ``Radio Silence persiste porque el miembro siente déficit de autonomía (percibe que sus preocupaciones no serán validadas) combinado con desconfianza en la benevolencia del equipo (cree que pueden usarlas en su contra)''
\item \textbf{Intervención calibrada:} En lugar de simplemente exhortar a ``comunicarse más'', la intervención se dirige a restaurar autonomía y benevolencia.
\end{itemize}

\item \textbf{Retrospective Deep-Dive Sessions:} Sesiones extendidas (1-2 horas, vs. las 1,5 horas típicas) cada 2-3 Sprints, dedicadas a inspeccionar profundamente un smell específico:
\begin{itemize}
\item Se invita a facilitador externo (coach ágil o especialista en dinámicas de grupo) si es necesario
\item Se usa technique como el ``World Café'' o ``Fishbowl'' para generar perspectives múltiples
\item El objetivo es entender no solo el síntoma, sino el contexto completo en que el smell prospera
\end{itemize}

\item \textbf{Continuous Pulse Check:} Surveys breves (2-3 preguntas) enviadas cada semana o cada dos semanas:
\begin{itemize}
\item ``En una escala de 1-5, ¿cuánto confías en las competencias de tu equipo?''
\item ``¿Hay algo que no hayas dicho esta semana pero debería ser dicho?''
\item ``¿Cómo fue tu conexión emocional con el equipo esta semana?''
\end{itemize}

Los resultados se grafican para detectar tendencias. Una caída repentina en confianza o un aumento en ``cosas no dichas'' es una señal de alerta que requiere inspección urgente.

\end{enumerate}

\subsubsection{Adaptación aplicada a la deuda social}

\paragraph{¿Por qué es crítica para la deuda social específicamente?}

La deuda social no se resuelve con una intervención única. Los community smells son síntomas de patrones profundos que requieren cambios sistémicos en cómo el equipo opera, comunica, y toma decisiones. La adaptación requiere que el equipo experimente, observe resultados, y ajuste continuamente.

Sin adaptación, la inspección se vuelve superficial: el equipo ve el problema, pero vuelve a los mismos patrones de comportamiento que lo generaron \cite{saeeda2024nontechnicaldebt}.

\paragraph{Ciclo de Adaptación para Deuda Social}

\begin{enumerate}

\item \textbf{Diseño de intervención social:} Basado en el análisis de causas (inspección), el equipo diseña una intervención específica:
\begin{itemize}
\item Objetivo claro: ``Reducir Radio Silence permitiendo que [miembro] se sienta seguro comunicando preocupaciones''
\item Mecanismo: ``Implementar estructura de feedback ``Feedback First'' en Daily Scrum, donde la primera pregunta es: ¿Hay algo que no estamos viendo?''
\item Hipótesis: ``Si creamos un ritual donde las preocupaciones son solicitadas explícitamente y acogidas sin defensa, Radio Silence disminuirá''
\item Duración: 2-3 Sprints para observar cambio
\item Dueño: Social Guide es responsable de facilitación; Scrum Master de sostener el ritual
\end{itemize}

\item \textbf{Implementación en el Sprint Backlog:} La intervención se incluye explícitamente en el Sprint Backlog Social:
\begin{itemize}
\item Tarea: ``Implementar Feedback First Protocol en Daily Scrum''
\item Estimación: 30 minutos de setup, 5 minutos adicionales por Daily
\item Done Criteria: Protocolo está documentado, equipo ha practicado 2 veces, todos entienden el flujo
\end{itemize}

\item \textbf{Ejecución con monitoreo:} Durante el Sprint, la intervención se ejecuta y se monitorea:
\begin{itemize}
\item Social Guide observa Daily Scrum y toma notas: ¿Alguien expresó una preocupación? ¿Cómo respondió el equipo?
\item Si la respuesta fue de cierre (``no, está bien''), el protocolo falló y requiere ajuste
\item Si la respuesta fue de apertura (``gracias por compartir, exploremos junto''), hay progreso
\end{itemize}

\item \textbf{Evaluación en Retrospective:} En la Sprint Retrospective Social:
\begin{itemize}
\item \textbf{Métricas:} ¿Radio Silence disminuyó? (medido por número de preocupaciones comunicadas vs. Sprints anteriores)
\item \textbf{Feedback cualitativo:} ¿El miembro se siente más seguro? ¿El equipo?
\item \textbf{Sostenibilidad:} ¿El protocolo es agotador o intuitivo?
\end{itemize}

\item \textbf{Decisión de adaptar:}
\begin{itemize}
\item \textbf{Mantener:} Si el protocolo funcionó y es sostenible, se mantiene como parte de la rutina del equipo
\item \textbf{Iterar:} Si funcionó parcialmente, se ajusta (ej.: hacer la pregunta en privado en lugar de en grupo, si el miembro aún siente exposición)
\item \textbf{Abandonar y probar otra:} Si no funcionó, se explora otra hipótesis de causa y otra intervención
\end{itemize}

\end{enumerate}

\paragraph{Mecanismos de adaptación específicos en SocialScrum}

\begin{enumerate}

\item \textbf{Social Experimentation Board:} Tabla visual en la que el equipo registra:
\begin{itemize}
\item Smell identificado
\item Hipótesis de causa (usando framework TAD/Confianza)
\item Intervención propuesta
\item Estado (diseñada, en implementación, evaluada)
\item Resultado (funcionó / parcial / no funcionó)
\item Próximo paso
\end{itemize}

Esta tabla mantiene el experimento social visible y responsable.

\item \textbf{Adaptation Retrospective (cada 4-6 Sprints):} Revisión extendida de qué intervenciones sociales han funcionado, cuáles no, y cuál es el patrón:
\begin{itemize}
\item ¿Las intervenciones que funcionan tienden a ser sobre autonomía, benevolencia, o integridad?
\item ¿Hay ciertos tipos de smells que son más resistentes a cambio?
\item ¿El equipo está desarrollando capacidad de auto-regulación, o sigue dependiendo del Social Guide para cada intervención?
\end{itemize}

El objetivo es que el equipo internalize el ciclo de inspección-adaptación y se vuelva cada vez más autónomo.

\item \textbf{Escalation Protocol:} Si una intervención no funciona después de 3 iteraciones, se escala:
\begin{itemize}
\item Se consulta con coach externo o especialista en dinámicas de grupo
\item Se cuestiona si la causa raíz realmente se entiende, o hay capas más profundas no identificadas
\item Se puede ajustar la intervención más radicalmente (ej.: cambio de pairing, cambio de ritmo de trabajo, etc.)
\end{itemize}

\end{enumerate}

